% Vertex-topology figure for cubic-operator sectors of App H §C.
%
% Four parity-even cubic operators generate four distinct linearised
% vertex topologies after one factor of F is replaced by the magnetic
% background F_bar:
%   (a)  R   x F^2  ->  h - A - F_bar    (lambda; 3-leg, Gertsenshtein channel)
%   (b)  T^2 x F    ->  T - T  with F_bar insertion (kappa; 2-pt mass shift)
%   (c)  R^2 x F    ->  h - h  with F_bar insertion (rho;   2-pt self-energy)
%   (d)  dT  x F^2  ->  T - A - F_bar    (mu; 3-leg torsion-photon channel)
%
% Leg conventions (matched to vertex_topology_cs.tex):
%   h        - coil-decorated line
%   A        - snake-decorated line
%   T        - dashed line
%   F_bar    - solid line ending in a cross (background insertion)
% Vertex dot carries the coupling-family letter (lambda, kappa, rho, mu).
% Parity-odd epsilon-twins are indicated by a small "+ epsilon" badge.
\documentclass[border=3pt,tikz,svgnames]{standalone}
\usepackage{stix}
\usepackage{tikz}
\usetikzlibrary{decorations.markings,decorations.pathmorphing,arrows.meta,calc}
\usepackage[outline]{contour}
\contourlength{0.5pt}
\tikzset{>={Stealth[length=1.2mm,width=1.2mm]}}

\colorlet{GravitonLine}{MediumPurple!50!SlateBlue}
\colorlet{PhotonLine}{Goldenrod!70!DarkOrange}
\colorlet{TorsionLine}{MediumAquamarine!60!Teal}
\colorlet{BgLine}{Sienna!80!Black}
\colorlet{VertexDot}{black!80!White}
\colorlet{PanelFrame}{Silver!50!White}
\colorlet{TagText}{black!75}

% Line styles
\tikzset{
  graviton/.style = {
    draw=GravitonLine, line width=0.7pt,
    decorate, decoration={coil, aspect=0.6, segment length=1.4mm, amplitude=0.55mm}
  },
  photon/.style = {
    draw=PhotonLine, line width=0.7pt,
    decorate, decoration={snake, segment length=1.4mm, amplitude=0.45mm}
  },
  torsion/.style = {
    draw=TorsionLine, line width=0.9pt, dash pattern={on 1.4mm off 0.9mm}
  },
  bgleg/.style = {
    draw=BgLine, line width=0.7pt
  },
  vertexdot/.style = {
    fill=VertexDot, draw=VertexDot, circle, inner sep=0pt, minimum size=2.0pt
  },
  bginsertion/.style = {
    fill=BgLine, draw=BgLine, circle, inner sep=0pt, minimum size=1.6pt
  },
  legend/.style = {font=\footnotesize, TagText}
}

\newcommand{\bgcross}[2][0.9mm]{%
  % cross marker at coordinate #2 with arm half-length #1
  \draw[BgLine, line width=0.8pt]
    ($(#2)+(45:#1)$)  -- ($(#2)+(225:#1)$)
    ($(#2)+(135:#1)$) -- ($(#2)+(-45:#1)$);
}

\newcommand{\paneltitle}[1]{%
  % Call inside a panel scope where coordinate (BC) marks the
  % bottom-centre anchor.
  % #1: sector mathy label, e.g. $\mathcal{R}\!\times\!\mathcal{F}^2$
  \node[font=\small, anchor=north] at ($(BC)+(0,-0.15)$) {#1};
}

\newcommand{\epsilonbadge}{%
  % Call inside a panel scope where coordinate (TR) marks the
  % top-right anchor. Badge sits inside the corner so leg labels
  % can extend out to the top-right without collision.
  \node[font=\scriptsize, TagText!85, anchor=north east, inner sep=2pt]
    at (TR) {{$+\,\varepsilon$}};
}

\begin{document}

\newlength{\figL}
\setlength{\figL}{0.95cm}

\begin{tikzpicture}[
    x=\figL, y=\figL,
    line cap=round, line join=round,
    font=\normalsize
]

% Panel grid: 2x2, each panel ~3.6 wide, 3.0 tall (in figL units).
% Panel centres:
\pgfmathsetmacro{\panW}{3.8}
\pgfmathsetmacro{\panH}{3.6}

\coordinate (Apos) at (0,            \panH);   % top-left  (a)
\coordinate (Bpos) at (\panW,        \panH);   % top-right (b)
\coordinate (Cpos) at (0,            0);       % bottom-left  (c)
\coordinate (Dpos) at (\panW,        0);       % bottom-right (d)

% Faint frames around each panel for visual structure
\foreach \P in {Apos, Bpos, Cpos, Dpos} {
  \draw[PanelFrame, very thin]
    ($(\P)+(-1.55,-1.45)$) rectangle ($(\P)+(1.55, 1.45)$);
}

% =========================================================
% Panel (a):  R x F^2 vertex  ->  h - A - F_bar  (lambda)
% =========================================================
\begin{scope}[shift={(Apos)}]
  \coordinate (V)  at (0,0);
  \coordinate (BC) at (0,-1.45);
  \coordinate (TR) at (1.42, 1.32);
  % three legs from the vertex
  \coordinate (Lh) at (135:1.15);  % graviton up-left
  \coordinate (LA) at (45:1.15);   % photon  up-right
  \coordinate (LF) at (-90:1.15);  % background straight down

  \draw[graviton] (V) -- (Lh);
  \draw[photon]   (V) -- (LA);
  \draw[bgleg]    (V) -- (LF);
  \bgcross{LF}

  \node[vertexdot] at (V) {};

  \node[GravitonLine, font=\small] at ($(Lh)+(-0.1,0.18)$) {$h$};
  \node[PhotonLine,   font=\small] at ($(LA)+(0.20,0.05)$) {$A$};
  \node[BgLine,       font=\small] at ($(LF)+(0.22,0)$)    {$\bar F$};

  \node[font=\small, TagText, anchor=north west, inner sep=2pt] at (-1.55, 1.45) {(a)};
  \epsilonbadge
  \paneltitle{$\mathcal{R}\!\times\!\mathcal{F}^{2}$}
\end{scope}

% =========================================================
% Panel (b):  T^2 x F vertex  ->  T - T  with F_bar insertion  (kappa)
% =========================================================
\begin{scope}[shift={(Bpos)}]
  \coordinate (BC) at (0,-1.45);
  \coordinate (TR) at (1.42, 1.32);
  % Two torsion legs joined by an F_bar insertion at the midpoint
  \coordinate (Tl) at (-1.15, 0.1);
  \coordinate (Tr) at ( 1.15, 0.1);
  \coordinate (Vins) at (0, 0.1);

  \draw[torsion] (Tl) -- (Vins);
  \draw[torsion] (Vins) -- (Tr);

  \node[bginsertion] at (Vins) {};
  \bgcross[1.0mm]{Vins}

  \node[BgLine, font=\small] at ($(Vins)+(0,-0.45)$) {$\bar F$};

  \node[TorsionLine, font=\small] at ($(Tl)+(-0.18,0.18)$) {$T$};
  \node[TorsionLine, font=\small] at ($(Tr)+( 0.18,0.18)$) {$T$};

  \node[font=\small, TagText, anchor=north west, inner sep=2pt] at (-1.55, 1.45) {(b)};
  \epsilonbadge
  \paneltitle{$\mathcal{T}^{2}\!\times\!\mathcal{F}$}
\end{scope}

% =========================================================
% Panel (c):  R^2 x F vertex  ->  h - h  with F_bar insertion  (rho)
% =========================================================
\begin{scope}[shift={(Cpos)}]
  \coordinate (BC) at (0,-1.45);
  \coordinate (TR) at (1.42, 1.32);
  \coordinate (Hl) at (-1.15, 0.1);
  \coordinate (Hr) at ( 1.15, 0.1);
  \coordinate (Vins) at (0, 0.1);

  \draw[graviton] (Hl) -- (Vins);
  \draw[graviton] (Vins) -- (Hr);

  \node[bginsertion] at (Vins) {};
  \bgcross[1.0mm]{Vins}

  \node[BgLine, font=\small] at ($(Vins)+(0,-0.45)$) {$\bar F$};

  \node[GravitonLine, font=\small] at ($(Hl)+(-0.18,0.18)$) {$h$};
  \node[GravitonLine, font=\small] at ($(Hr)+( 0.18,0.18)$) {$h$};

  \node[font=\small, TagText, anchor=north west, inner sep=2pt] at (-1.55, 1.45) {(c)};
  \paneltitle{$\mathcal{R}^{2}\!\times\!\mathcal{F}$}
\end{scope}

% =========================================================
% Panel (d):  dT x F^2 vertex  ->  T - A - F_bar  (mu)
% =========================================================
\begin{scope}[shift={(Dpos)}]
  \coordinate (V)  at (0,0);
  \coordinate (BC) at (0,-1.45);
  \coordinate (TR) at (1.42, 1.32);
  \coordinate (LT) at (135:1.15);  % torsion up-left
  \coordinate (LA) at (45:1.15);   % photon  up-right
  \coordinate (LF) at (-90:1.15);  % background straight down

  \draw[torsion] (V) -- (LT);
  \draw[photon]  (V) -- (LA);
  \draw[bgleg]   (V) -- (LF);
  \bgcross{LF}

  \node[vertexdot] at (V) {};

  \node[TorsionLine, font=\small] at ($(LT)+(-0.1,0.18)$) {$T$};
  \node[PhotonLine,  font=\small] at ($(LA)+(0.20,0.05)$) {$A$};
  \node[BgLine,      font=\small] at ($(LF)+(0.22,0)$)    {$\bar F$};

  \node[font=\small, TagText, anchor=north west, inner sep=2pt] at (-1.55, 1.45) {(d)};
  \epsilonbadge
  \paneltitle{$\widetilde{\nabla}\mathcal{T}\!\times\!\mathcal{F}^{2}$}
\end{scope}

\end{tikzpicture}

\end{document}
